\subsection{HTTP, JSON và CORS}

\boncau
{Cách trình duyệt nói chuyện với máy chủ: mỗi lần gọi độc lập, thân REST là
JSON, khi dev hai cổng khác nhau thì trình duyệt hỏi CORS~\cite{rfc9110}.}
{Toàn bộ nghiệp vụ đi qua HTTP. Body REST là JSON. Khi Vite (5173) gọi API
(8023) khác origin, trình duyệt bắt buộc có CORS.}
{GET đọc, POST tạo, PATCH sửa, DELETE xóa. Mã dùng: 200, 201, 202 (soát xét
nền), 400, 401, 403, 404 (kể cả xem chéo hội thoại), 422, 429, 503. CORS khóa
origin cấu hình, không \texttt{*}. Docker cùng origin nên CORS gần như không chạy.}
{Mọi request API. OPTIONS chỉ khi trình duyệt preflight. HTTPS/TLS không bật
trong Compose nội bộ --- khi đưa ra Internet phải đặt trước nginx.}

\subsection{Python}

\boncau
{Ngôn ngữ của toàn bộ tầng máy chủ: FastAPI, tách từ, LangGraph và SQLAlchemy
chung một codebase.}
{\texttt{underthesea} (tách từ BM25) ghim Python 3.10; FastAPI, SQLAlchemy,
LangGraph cùng một ngôn ngữ.}
{Quản lý gói bằng \texttt{uv}; type hint + Pydantic cho schema; \texttt{async}
cho I/O tới Postgres, Chroma và Gemini.}
{Toàn bộ tầng máy chủ. Không dùng cho frontend.}

\subsection{FastAPI}

\boncau
{Cửa HTTP duy nhất của SPA: nhận REST và đẩy SSE, kiểm body bằng Pydantic, tự
sinh mô tả OpenAPI~\cite{fastapi}.}
{Một request hỏi đáp chủ yếu chờ mạng: bất đồng bộ có lợi. SSE có
\texttt{StreamingResponse}. Schema khớp TypeScript phía trình duyệt.}
{Router tách theo nghiệp vụ (\texttt{auth}, hội thoại, hỏi đáp, luật, tài liệu,
lịch, quản trị). Pipeline RAG nằm tầng dịch vụ, không nằm trong router.}
{Mọi HTTP nghiệp vụ. Đường \texttt{/docs} dùng khi trình bày API, không phải
luồng người dùng DNNVV.}

\subsection{Pydantic}

\boncau
{Lớp biên giới API: body sai thì 422, không vào nghiệp vụ; schema khớp
TypeScript phía trình duyệt.}
{Lỗi 422 có cấu trúc khi client gửi sai trường; OpenAPI khớp schema thật.}
{Body đăng ký, đăng nhập, tin nhắn, kết quả soát xét, bộ câu chạy hàng loạt.}
{Mỗi lần FastAPI nhận hoặc trả JSON có khai báo mô hình.}

\subsection{SQLAlchemy 2.0 và Alembic}

\boncau
{Cách đọc ghi PostgreSQL bằng lớp Python, kèm nhật ký đổi cột để máy khác lên
cùng schema.}
{CRUD hội thoại, tài liệu, lịch không viết SQL lặp. Không sửa schema tay trên
máy rồi quên ghi.}
{Bản bất đồng bộ + \texttt{asyncpg}. User ứng dụng không \texttt{CREATE EXTENSION};
extension nằm \texttt{init.sql}.}
{Mọi đọc/ghi PostgreSQL nghiệp vụ. Vector không đi ORM --- nằm Chroma.}

\subsection{PostgreSQL}

\boncau
{Nơi giữ user, hội thoại, Điều chữ, lịch và ràng buộc (mã số thuế duy nhất, task
không trùng kỳ).}
{Cần ràng buộc tổ chức--user--hội thoại, FTS tiếng Việt tự dựng, và lịch tuân
thủ. Chroma không làm việc này.}
{Bản 17. Cổng host 23432. Corpus chữ và nghiệp vụ cùng một cụm; cột generated
\texttt{tsvector} sau \texttt{unaccent}.}
{Luôn chạy khi hệ thống sống. Tra cứu FTS không cần Gemini.}

\subsection{ChromaDB}

\boncau
{Nơi giữ vector từng Điều và tìm $k$ Điều gần câu hỏi theo cosine, chạy trong
cùng tiến trình backend.}
{Nhẹ, chạy kèm tiến trình backend, đủ \soDieu{} Điều; không thêm cụm Pinecone.}
{Volume \texttt{chroma\_db}. Manifest (mô hình, chiều, số bản ghi) so khớp lúc
khởi động. Đổi chiều thì xóa và nhúng lại.}
{Nạp/reindex và nút \texttt{retrieve}. Thiếu chỉ mục: hỏi đáp 503, FTS vẫn dùng.}

\subsection{LangGraph}

\boncau
{Cách nối bảy bước hỏi đáp thành một đồ thị có nhánh thoát sớm và có sự kiện
SSE từng nút.}
{Cần thoát sớm khi câu không phải pháp luật, và đẩy SSE theo từng nút. Không
dùng agent tự chọn công cụ tự do: thứ tự truy hồi rồi sinh là cố định.}
{Bảy nút; bộ nhớ ngắn nạp từ PostgreSQL, không dùng \texttt{InMemorySaver}.}
{Mỗi lượt hỏi đáp. Soát xét hợp đồng tái sử dụng hàm truy hồi, không đi đủ bảy
nút hội thoại.}

\subsection{Uvicorn}

\boncau
{Tiến trình chạy FastAPI trong Docker: một worker, cổng 8023, đủ để giữ SSE và
không đua file Chroma.}
{SSE cần vòng sự kiện bất đồng bộ; Gunicorn đồng bộ thuần không hợp.}
{Docker: một tiến trình, không reload, cổng 8023. Một worker để hai process
không đua file Chroma.}
{Luôn khi backend sống. Reload chỉ trên máy phát triển ngoài Compose.}

\subsection{Mô hình ngôn ngữ và Gemini}

\boncau
{Dịch vụ viết câu và nhúng chữ qua API (Gemini), không phải kho luật và không
được tinh chỉnh trọng số trong đồ án.}
{Không có GPU để tự host. Lớp tương thích OpenAI cho phép đổi \texttt{base\_url}
sang vLLM sau này mà không viết lại nút.}
{\texttt{gemini-2.5-flash}, nhiệt độ 0, tối đa 4096 token. Nhúng tách
\texttt{gemini-embedding-001}. Flash thay Pro: đủ câu ngắn trên $k=8$ Điều, rẻ
hơn trên HyDE + filter + generate.}
{Intent, HyDE, filter, generate, và nhúng. Thiếu khóa: health 200, chat 503.
Không fine-tune.}

\subsection{Underthesea}

\boncau
{Công cụ tách từ tiếng Việt cho BM25, để ``lao động'' không bị đếm thành hai
mảnh.}
{BM25 cần đơn vị từ. ``Lao động'' tách thành ``lao'' / ``động'' làm nhiễu IDF.}
{Nạp được thì dùng; không thì fallback regex. Cache theo dấu vân tay tokenizer.}
{Lúc dựng chỉ mục BM25 (khởi động hoặc cache lệch). Không dùng cho embedding.}

\subsection{React}

\boncau
{Thư viện vẽ toàn bộ màn hình sau đăng nhập; đổi trang không tải lại cả
site~\cite{react}.}
{Đề bài yêu cầu React. Ứng dụng sau đăng nhập, không cần SSR/SEO.}
{React 19. Tách \texttt{ChatPanel}, \texttt{CitationList}, \texttt{AgentTrace}.
Không dùng Next.js.}
{Mọi màn hình sau login. Trang login/register cũng là component React.}

\subsection{TypeScript}

\boncau
{JavaScript có kiểu: lệch tên trường API (token, citation) bị bắt lúc build.}
{Lệch \texttt{access\_token} / \texttt{accessToken} bắt lúc biên dịch.}
{Khai báo user, hội thoại, Điều, finding, sự kiện SSE. Khớp OpenAPI FastAPI.}
{Toàn bộ frontend.}

\subsection{Vite}

\boncau
{Công cụ dựng SPA: lúc lập trình có proxy \texttt{/api}, lúc đóng gói ra tệp
tĩnh cho nginx.}
{Khởi động nhanh; proxy \texttt{/api} giống nginx production.}
{\texttt{VITE\_API\_BASE\_URL} trống trong Docker: cùng origin.}
{Dev trên máy; image frontend chỉ chứa bản đã build.}

\subsection{Tailwind CSS}

\boncau
{Cách trang trí giao diện bằng class nhỏ trong JSX, không kéo bộ nút thương mại.}
{Không kéo bộ component thương mại; gói nhỏ, nền tối đủ đọc Điều dài.}
{Bản 4. Nút, ô nhập, thẻ tự viết.}
{Mọi trang giao diện.}

\subsection{TanStack Query}

\boncau
{Bộ nhớ tạm phía trình duyệt cho danh mục luật và trạng thái soát xét, tách
khỏi kho phiên đăng nhập.}
{Danh mục luật, cây Điều, poll soát xét --- không nhét vào store phiên.}
{Poll 2,5 giây khi soát xét. Tắt \texttt{refetchOnWindowFocus}.}
{Mọi GET danh sách/chi tiết. Không dùng cho luồng SSE (Fetch).}

\subsection{Zustand}

\boncau
{Kho nhỏ phía trình duyệt: chỉ user, token và hàm khôi phục phiên.}
{Chỉ cần user + token + \texttt{restore}; Redux thừa boilerplate.}
{Đọc token \texttt{localStorage} khi tải trang. Không chứa danh sách luật.}
{Khởi động SPA và mỗi lần login/logout.}

\subsection{React Router}

\boncau
{Cách đổi URL mà không tải lại trang; số hiệu có dấu \texttt{/} cần bắt phần
còn lại của đường dẫn.}
{Số hiệu văn bản chứa \texttt{/} cần splat \texttt{*}.}
{Route \texttt{/login}, \texttt{/chat/:id}, \texttt{/laws/:lawId/*},
\texttt{/contracts}, \texttt{/compliance}, \texttt{/admin/*},
\texttt{/admin/batch-eval}. Admin \texttt{lazy()}. \texttt{RequireAuth} chờ
restore xong.}
{Mỗi lần người dùng đổi trang.}

\subsection{Axios}

\boncau
{Client REST: gắn Bearer, làm mới token một lần khi 401, gỡ JSON khi gửi file.}
{Cần interceptor chung: Bearer, làm mới token một lần, hàng đợi 401.}
{Xóa \texttt{Content-Type} khi \texttt{FormData} để trình duyệt đặt multipart.}
{Mọi REST. Không dùng cho SSE --- hỏi đáp dùng Fetch.}

\subsection{nginx}

\boncau
{Cổng duy nhất người dùng mở: phát tệp Vite và chuyển \texttt{/api} vào
backend, tắt đệm cho SSE.}
{Trình duyệt chỉ mở một cổng; \texttt{/api} vào backend.}
{Image frontend phục vụ bản Vite. Host 5173.}
{Môi trường Docker / production. Dev dùng proxy Vite thay nginx.}

\subsection{Docker và Docker Compose}

\boncau
{Cách dựng cùng một cụm Postgres--API--nginx trên máy khác nhau, không lệch
phiên bản Python hay CSDL~\cite{docker}.}
{Cùng Python 3.10, Postgres 17, nginx trên máy khác nhau. Không lệch phiên bản
khi chấm đồ án.}
{Ba dịch vụ: \texttt{postgres} (healthcheck \texttt{pg\_isready}),
\texttt{backend} (đợi healthy, \texttt{/health}), \texttt{frontend}. Compose
chỉ đọc \texttt{.env} gốc, không đọc \texttt{backend/.env}.}
{Lúc dựng và vận hành cụm. Máy dev có thể chạy Vite + Postgres ngoài Compose.}

\subsection{JWT}

\boncau
{Hai loại token ký sẵn (access ngắn, refresh dài) để API biết ai đang gọi mà
không giữ phiên trên RAM~\cite{rfc7519}.}
{Không thêm Redis. Access ngắn hạn giảm cửa sổ lộ token.}
{HS256, \texttt{JWT\_SECRET\_KEY}. Claim \texttt{sub}, \texttt{type}
(\texttt{access}/\texttt{refresh}), \texttt{iat}, \texttt{exp}, \texttt{jti}.
Access 15 phút, refresh 7 ngày. Gọi API bằng refresh bị từ chối.}
{Mọi request sau login. Token ở \texttt{localStorage} --- còn rủi ro XSS; hướng
phát triển là cookie \texttt{HttpOnly}.}

\subsection{Argon2}

\boncau
{Cách lưu mật khẩu: băm argon2id, cố ý chậm, không lưu chữ gốc~\cite{argon2}.}
{SHA-256 quá nhanh, dễ dò GPU. Cần hàm cố ý chậm và tốn bộ nhớ.}
{Qua passlib. Không lưu plaintext. Hash hỏng không sập login --- verify trả sai.}
{Đăng ký và mỗi lần đăng nhập. Không dùng cho JWT.}

\subsection{Server-Sent Events}

\boncau
{Kênh máy chủ đẩy tiến độ và từng chữ xuống trình duyệt trong một kết nối HTTP,
không mở WebSocket.}
{Cần tiến độ từng nút và từng token. WebSocket thừa chiều client; polling chat
giật chữ.}
{Ba loại: \texttt{status}, \texttt{token}, \texttt{result}. Heartbeat tránh
proxy cắt. Client dùng Fetch (POST + Bearer) vì \texttt{EventSource} không gắn
header.}
{Mỗi lượt hỏi đáp. Soát xét hợp đồng dùng polling 2,5 giây, không SSE.}

\subsection{slowapi}

\boncau
{Chặn một tài khoản gọi chat hoặc tải file quá dày, để khỏi vét hạn mức mô hình.}
{Một tài khoản không được vét hạn mức Gemini.}
{Chat 20 lần/phút, tải tệp 30 lần/giờ, theo user (hoặc IP nếu chưa login).}
{Các endpoint tốn mô hình hoặc tốn đĩa. Tra cứu FTS không siết cùng mức chat.}

\subsection{Kiến trúc client--server và SPA}

\boncau
{Tách giao diện và máy chủ: trình duyệt tải một bộ JS, mọi nghiệp vụ đi
\texttt{/api/v1}.}
{Tách giao diện và máy chủ; trang sau login không cần HTML sẵn cho SEO.}
{Tiền tố \texttt{/api/v1}. Ba môi trường: Vite + Postgres local, Compose ba
container, tương lai tách domain.}
{Toàn bộ sản phẩm. Không dùng SSR.}
